\documentclass[11pt]{article}

\usepackage[margin=1in]{geometry}
\usepackage{amsmath,amssymb,amsthm,mathtools}
\usepackage{enumitem}
\usepackage[hidelinks]{hyperref}
\usepackage{microtype}

\allowdisplaybreaks
\setlist[itemize]{topsep=2pt,itemsep=1pt,leftmargin=1.5em}
\setlist[enumerate]{topsep=2pt,itemsep=1pt,leftmargin=1.6em}

\newtheorem{theorem}{Theorem}[section]
\newtheorem{proposition}[theorem]{Proposition}
\newtheorem{lemma}[theorem]{Lemma}
\newtheorem{corollary}[theorem]{Corollary}
\theoremstyle{definition}
\newtheorem{definition}[theorem]{Definition}
\theoremstyle{remark}
\newtheorem{remark}[theorem]{Remark}

\newcommand{\nuu}{\nu_2}
\newcommand{\Ecal}{\mathcal E}
\newcommand{\Rcal}{\mathcal R}
\newcommand{\V}{\mathcal V}
\newcommand{\GD}{G_D}
\newcommand{\CD}{C_D}
\newcommand{\bdD}{\partial_D}
\newcommand{\depth}{\operatorname{depth}}
\newcommand{\dist}{\operatorname{dist}}
\newcommand{\oddrep}{\operatorname{oddrep}}
\newcommand{\Pre}{\mathsf{Pre}}
\newcommand{\Seed}{\operatorname{Seed}}
\newcommand{\Pair}{\mathsf P}
\newcommand{\Sfam}{\mathcal S}
\newcommand{\eps}{\varepsilon}

\title{Adaptive Valuation Graphs and Binomial Distance Layers\\
\large A finite graph-theoretic structure from the \texorpdfstring{$3n+1$}{3n+1} valuation rule}
\author{}
\date{Second draft --- April 2026}

\begin{document}
\maketitle

\begin{abstract}
For each integer $D\ge 1$ we define a finite directed graph $G_D$ on odd residue classes modulo powers of two up to depth $D$.  The graph is adaptive: a class whose two-adic valuation data do not yet determine a positive-depth transition refines by one binary digit, while every other class carries the valuation-labeled transition determined by that data.  We prove that $G_D$ has a unique nontrivial strongly connected component $C_D$, give an explicit description of it, and show that $|C_D|=D(D-2)$ for $D\ge3$.  We also prove that valuation-edge preimages of a target of depth $e$ are naturally indexed by pairs $e<i<j\le D$.  This pair indexing yields a canonical history coordinate for vertices at positive distance from $C_D$.  After quotienting by the top-bit involution, the depth-$d$, distance-$k$ layer is naturally bijective with
\[
\Sfam_{d,k}=\{\Sigma\subseteq[d]: d\in\Sigma,\ |\Sigma|=2k+2\},
\]
and hence the oriented layer has cardinality
\[
\bigl|\{v:\depth(v)=d,\ R_D(v)=k\}\bigr|=2\binom{d-1}{2k+1}\qquad(k\ge1).
\]
The Collatz map motivates the local valuation rule, but the results are finite graph theory and enumerative combinatorics.
\end{abstract}

\section{Introduction}

This paper studies a finite directed graph built from odd residue classes modulo powers of two.  The graph is not drawn at one fixed modulus.  Instead, it refines a residue class only when the currently visible two-adic valuation data are insufficient to define a positive-depth transition.  This adaptive rule produces a finite graph with a rigid structure: a single recurrent core, a small truncation boundary, pair-indexed valuation preimages, and binomial distance layers.

The construction comes from the following two-adic valuation rule.  For an odd integer $X$, write
\[
X+1=2^tq,\qquad q\text{ odd},
\]
and then
\[
3^tq-1=2^\rho Y,\qquad Y\text{ odd}.
\]
The finite graph asks what can be determined from the residue class of $X$ modulo $2^d$.  If the class already determines the valuation data $(t,\rho)$ up to a positive target depth, the graph follows the induced transition.  If not, it asks for one more binary digit.

The main point is that the theorems below are not integer-orbit statements.  The graph and its combinatorics are the subject.  Collatz is the source of the valuation rule.

\section{The graph}\label{sec:graph}

For $d\ge1$, let $V_d$ be the set of odd residues modulo $2^d$, represented by odd integers $1\le r<2^d$.  A vertex of $G_D$ is a pair $(r,d)$ with $1\le d\le D$ and $r\in V_d$; write $\depth(r,d)=d$.  Thus
\[
V(G_D)=\{(r,d):1\le d\le D,\ r\equiv1\pmod2,\ 1\le r<2^d\}.
\]
If $x$ is odd modulo $2^m$, let $\oddrep_m(x)$ denote its canonical representative in $\{1,3,\ldots,2^m-1\}$.

\subsection{Exceptional residues and valuation edges}

For $1\le t<d$, define
\[
r_{d,t}\equiv 2^t(3^t)^{-1}-1\pmod {2^d},
\]
where the inverse is taken modulo $2^{d-t}$ and the odd representative modulo $2^d$ is chosen.  Also set
\[
r_{d,d}=2^d-1,
\qquad
\Ecal_d=\{r_{d,1},\ldots,r_{d,d}\}.
\]
The residues in $\Ecal_d$ are distinct, since $\nuu(r_{d,t}+1)=t$ for $t<d$ and $\nuu(r_{d,d}+1)=d$.  We call them exceptional.  All other residues at depth $d$ are regular.

The edges of $G_D$ are as follows.

If $(r,d)$ is exceptional and $d<D$, draw the two refinement edges
\[
(r,d)\to(r,d+1),\qquad (r,d)\to(r+2^d,d+1).
\]
If $d=D$, no refinement edge is drawn.

If $(r,d)$ is regular, put
\[
t=\nuu(r+1),\qquad q=\frac{r+1}{2^t},\qquad
\rho=\nuu(3^tq-1),\qquad e=d-t-\rho.
\]
Lemma~\ref{lem:local} below implies $e\ge1$.  Set
\[
y=\frac{3^tq-1}{2^\rho},\qquad s=\oddrep_e(y),
\]
and draw the valuation edge
\[
(r,d)\to(s,e).
\]
We denote this partial map by $F(r,d)=(s,e)$.

The truncation boundary is
\[
\bdD=\{(r,D):r\in\Ecal_D\}\cup\{(2^{D-1}-1,D-1)\}.
\]
The second term is the maximal exceptional vertex one level below the cutoff.

\begin{lemma}[Local valuation algebra]\label{lem:local}
The following hold.
\begin{enumerate}
\item A regular vertex has a valid valuation edge: $e=d-t-\rho\ge1$.
\item Every exceptional vertex at depth $d>1$ reduces to an exceptional vertex at depth $d-1$.
\item If $1\le t<d$, then the two children of $(r_{d,t},d)$ consist of one exceptional child of type $t$ and one regular child whose valuation edge has target $(1,1)$.
\item The two children of the maximal exceptional vertex $(2^d-1,d)$ are exceptional of types $d$ and $d+1$ at depth $d+1$.
\end{enumerate}
\end{lemma}

\begin{proof}
For a regular vertex, $e\le0$ would mean
\[
3^tq\equiv1\pmod {2^{d-t}},
\]
so $r=2^tq-1$ would be the exceptional residue of type $t$, a contradiction.

The parent statement follows immediately from the congruence defining $r_{d,t}$; the case $t=d$ reduces to the maximal exceptional residue at depth $d-1$.

For the child statement, write $r=r_{d,t}$ with $t<d$ and $q=(r+1)/2^t$.  The two depth-$(d+1)$ lifts replace $q$ by $q$ and $q+2^{d-t}$.  Since
\[
3^t(q+2^{d-t})-1=(3^tq-1)+3^t2^{d-t},
\]
exactly one lift satisfies the stronger congruence modulo $2^{d+1-t}$; this is the exceptional child of type $t$.  The other has exact valuation $\rho=d-t$, hence target depth $(d+1)-t-(d-t)=1$, so it maps to $(1,1)$.  Finally, the children of $2^d-1$ are $2^d-1$ and $2^{d+1}-1$, which are the exceptional residues of types $d$ and $d+1$ at depth $d+1$.
\end{proof}

A regular vertex is a \emph{direct return} if its valuation edge lands at $(1,1)$.  Let $\Rcal_d$ be the set of direct returns at depth $d$.  Lemma~\ref{lem:local} gives
\[
|\Rcal_d|=\max(d-2,0),
\]
and for $d\ge3$ these are exactly the regular children of the non-maximal exceptional vertices at depth $d-1$.

\section{The recurrent core}\label{sec:core}

For $D\ge3$, define
\[
K_D=\left(\{(r,d):r\in\Ecal_d,\ 1\le d\le D\}\cup\bigcup_{d=1}^D\Rcal_d\right)\setminus\bdD.
\]
Thus $K_D$ consists of the exceptional vertices and direct returns, except for the truncation boundary.

\begin{theorem}[Core theorem]\label{thm:core}
For $D\ge3$, $G_D$ has a unique nontrivial strongly connected component $C_D$, and
\[
C_D=K_D.
\]
Moreover
\[
|C_D|=D(D-2),
\]
and a vertex reaches $C_D$ if and only if it is not in $\bdD$.
\end{theorem}

\begin{proof}
First, $K_D$ is strongly connected.  Every direct return maps to $(1,1)$.  If an exceptional vertex in $K_D$ has non-maximal type, Lemma~\ref{lem:local} gives a direct-return child and hence a path to $(1,1)$.  If it is maximal, it is not in the boundary, so its depth is at most $D-2$; one refinement step reaches a non-maximal exceptional child, and then the previous case applies.  Conversely, every exceptional vertex is reachable from $(1,1)$ by repeatedly taking exceptional parents, and every direct return is a refinement child of an exceptional vertex.  Hence $(1,1)$ reaches all of $K_D$.

Now let $v\notin K_D\cup\bdD$.  Then $v$ is regular and not a direct return.  Its unique valuation edge strictly lowers depth, since $t\ge1$ and $\rho\ge1$.  Iterating therefore reaches either $K_D$ or the boundary.  It cannot first hit the boundary: every valuation edge from a source of depth at most $D$ has target depth at most $D-2$, while the boundary lies only in depths $D-1$ and $D$.  Hence every non-boundary vertex reaches $K_D$.

The boundary does not reach $K_D$: terminal exceptional vertices have no outgoing refinement edges, and the extra boundary vertex $(2^{D-1}-1,D-1)$ refines only to terminal exceptional vertices.  Finally, outside $K_D$ any path either enters the boundary trap or follows valuation edges of strictly decreasing depth, so no further directed cycle is possible.  Thus $K_D$ is the unique nontrivial strongly connected component.

For the count,
\[
\sum_{d=1}^D |\Ecal_d|=\sum_{d=1}^D d=\frac{D(D+1)}2,
\]
and
\[
\sum_{d=1}^D |\Rcal_d|=\sum_{d=3}^D(d-2)=\frac{(D-1)(D-2)}2.
\]
Their union has size $D^2-D+1$.  Removing the $D$ terminal exceptional vertices and the additional boundary vertex gives
\[
|C_D|=D^2-D+1-(D+1)=D(D-2).
\]
\end{proof}

Let
\[
R_D(v)=\dist(v,C_D)
\]
when $v$ reaches $C_D$, and $R_D(v)=\infty$ otherwise.  Theorem~\ref{thm:core} says that $R_D(v)<\infty$ exactly off the boundary.  If $R_D(v)=k>0$, then $v$ is regular and not a direct return, and repeated valuation edges give a unique chain
\[
v=v_k\to v_{k-1}\to\cdots\to v_0\in C_D.
\]

\section{Preimages and pair indexing}\label{sec:preimages}

The next theorem is the basic enumerative mechanism.  It is target-first: fix a target and solve for all valuation-edge sources.

\begin{theorem}[Preimage-pair theorem]\label{thm:preimage}
Fix a target vertex $u=(s,e)\in V(G_D)$.  The valuation-edge preimages of $u$ are naturally indexed by pairs
\[
e<i<j\le D.
\]
For the pair $(i,j)$, set
\[
t=i-e,
\qquad
\rho=j-i,
\qquad
 d=j.
\]
Then the unique source is obtained by solving
\[
q\equiv3^{-t}(1+2^\rho s)\pmod {2^{e+\rho}}
\]
and setting
\[
\Pre_{i,j}(s,e)=\bigl(\oddrep_j(2^tq-1),j\bigr).
\]
Consequently $|F^{-1}(s,e)|=\binom{D-e}{2}$.
\end{theorem}

\begin{proof}
If a valuation edge has source depth $d$ and target depth $e$, then
\[
d=e+t+\rho.
\]
Thus $i=e+t$ and $j=d$ satisfy $e<i<j\le D$.

Conversely, fix $e<i<j\le D$ and put $t=i-e$, $\rho=j-i$.  A source with target residue $s$ must satisfy
\[
3^tq-1\equiv2^\rho s\pmod {2^{e+\rho}},
\]
or equivalently
\[
q\equiv3^{-t}(1+2^\rho s)\pmod {2^{e+\rho}}.
\]
This congruence has a unique odd solution modulo $2^{e+\rho}$.  With $r=\oddrep_j(2^tq-1)$, one has $\nuu(r+1)=t$.  Since $s$ is odd, the valuation of $3^tq-1$ is exactly $\rho$, and division by $2^\rho$ gives target residue $s$ modulo $2^e$.  The source is regular because $e\ge1$.  Uniqueness follows from the same congruence for $q$.
\end{proof}

The pair attached to a valuation edge $v=(r,d)\to(s,e)$ is therefore
\[
\Pair(v)=\{e+\nuu(r+1),d\}.
\]
Each preimage step adds two depths: an intermediate valuation depth and the source depth.

\section{Histories and binomial layers}\label{sec:histories}

For $d\ge2$, let $\tau_d$ be the top-bit involution
\[
\tau_d(r,d)=
\begin{cases}
(r+2^{d-1},d),&r<2^{d-1},\\
(r-2^{d-1},d),&r>2^{d-1}.
\end{cases}
\]
It exchanges the two depth-$d$ lifts of the same residue modulo $2^{d-1}$.

\begin{proposition}[Stable core seeds]\label{prop:seeds}
For $2\le d\le D-2$, the quotient $(C_D\cap V_d)/\tau_d$ is canonically labeled by the two-element sets
\[
\{a,d\},\qquad 1\le a<d.
\]
There are two oriented core vertices above each label.
\end{proposition}

\begin{proof}
For $1\le a\le d-2$, the two children of the exceptional vertex of type $a$ at depth $d-1$ consist of the exceptional vertex of type $a$ at depth $d$ and its direct-return sibling.  These two children differ by $2^{d-1}$, so they form one $\tau_d$-orbit.  The remaining orbit comes from the two children of the maximal exceptional vertex at depth $d-1$, namely the exceptional vertices of types $d-1$ and $d$.  Since $d\le D-2$, no boundary vertex is involved.  These $d-1$ two-element orbits exhaust $C_D\cap V_d$.
\end{proof}

If $v=v_k\to\cdots\to v_0\in C_D$ is the unique chain for a vertex of positive finite distance, write $v_m=(r_m,d_m)$.  The endpoint $v_0$ has stable depth: $2\le d_0\le D-2$, because the final edge is not a direct return and every valuation edge lowers depth by at least two.  Let
\[
\Seed(v_0)=\{a,d_0\}
\]
be its core seed label from Proposition~\ref{prop:seeds}.  Define the projected history
\[
\Sigma(v)=\Seed(v_0)\cup\bigcup_{m=1}^k\{d_{m-1}+\nuu(r_m+1),d_m\}.
\]
For a stable core vertex itself, set $\Sigma(v_0)=\Seed(v_0)$.

\begin{lemma}[History intervals]\label{lem:historyintervals}
If $v$ has finite distance $k$ and depth $d$, then the union defining $\Sigma(v)$ is disjoint.  For $k\ge1$,
\[
\Sigma(v)\subseteq[d],\qquad d\in\Sigma(v),\qquad |\Sigma(v)|=2k+2.
\]
\end{lemma}

\begin{proof}
For the edge $v_m\to v_{m-1}$, put
\[
t_m=\nuu(r_m+1),\qquad \rho_m=\nuu(3^{t_m}q_m-1).
\]
Then
\[
d_m=d_{m-1}+t_m+\rho_m,
\]
so
\[
\{d_{m-1}+t_m,d_m\}\subset(d_{m-1},d_m].
\]
The seed lies in $[1,d_0]$.  Inducting along the chain, each preimage step appends two new elements in the next interval $(d_{m-1},d_m]$.  Hence the union is disjoint, lies in $[d]$, contains the final depth $d$, and has cardinality $2+2k$.
\end{proof}

For $k\ge1$, define
\[
\Sfam_{d,k}=\{\Sigma\subseteq[d]:d\in\Sigma,\ |\Sigma|=2k+2\}.
\]

\begin{theorem}[History bijection]\label{thm:bijection}
For $D\ge3$, $k\ge1$, and $1\le d\le D$, the map
\[
v\longmapsto (\eps(v),\Sigma(v))
\]
is a natural bijection
\[
\{v:\depth(v)=d,\ R_D(v)=k\}
\longrightarrow
\{+,-\}\times\Sfam_{d,k},
\]
where $\eps(v)$ is the orientation of the terminal core seed.  Consequently
\[
\bigl|\{v:\depth(v)=d,\ R_D(v)=k\}\bigr|=2\binom{d-1}{2k+1}.
\]
\end{theorem}

\begin{proof}
Lemma~\ref{lem:historyintervals} shows that the map lands in $\{+,-\}\times\Sfam_{d,k}$.

To reconstruct, take $(\eps,\Sigma)$ with
\[
\Sigma=\{\sigma_1<\sigma_2<\cdots<\sigma_{2k+2}=d\}.
\]
Start with the oriented core vertex labeled by $\{\sigma_1,\sigma_2\}$ and orientation $\eps$.  At step $m=1,\ldots,k$, use Theorem~\ref{thm:preimage} to take the unique valuation-edge preimage of the current target of depth $\sigma_{2m}$ indexed by the pair
\[
\sigma_{2m}<\sigma_{2m+1}<\sigma_{2m+2}.
\]
This constructs a chain whose final source has depth $d$.

Each reconstructed edge has target depth at least $2$, so no reconstructed source is a direct return; hence the first time the chain reaches the core is its prescribed endpoint.  Thus the final source has distance $k$.  The construction is inverse to the history map because Theorem~\ref{thm:preimage} gives a unique source for each target and pair.

The count follows from
\[
|\Sfam_{d,k}|=\binom{d-1}{2k+1},
\]
since $d$ is forced to lie in $\Sigma$.
\end{proof}

The top-bit quotient removes the two orientations.  The needed equivariance is elementary: if a valuation edge has target depth $e\ge2$, then replacing the source by its top-bit flip changes $q$ to $q\pm2^{d-1-t}$, preserves the exact valuations $(t,\rho)$, and flips the top bit of the target.  Thus valuation chains commute with $\tau$ as long as no direct return is involved.  Positive-distance chains have this property.

\begin{corollary}[Quotient layers]\label{cor:quotient}
For $k\ge1$ and $d\ge2$,
\[
\{v:\depth(v)=d,\ R_D(v)=k\}/\tau_d
\cong
\Sfam_{d,k}.
\]
Equivalently, summing over $d\le D$,
\[
\bigl|\{v:R_D(v)=k\}\bigr|=2\binom{D}{2k+2}.
\]
\end{corollary}

\section{Relation to Collatz}\label{sec:collatz}

The rule used to define $G_D$ is extracted from a residue-class model of the odd Collatz/Syracuse dynamics.  In that model one studies the odd-to-odd valuation update
\[
X+1=2^tq,
\qquad
3^tq-1=2^\rho Y.
\]
This origin explains why the formulas occur, but it is not needed for the finite graph theorems.

The statements proved here do not assert that integer Collatz orbits reach $1$.  A vertex $(r,d)$ is a residue class, not an integer.  The integer $1$ is represented by an infinite compatible tower of residue classes, and integer dynamics require lift data beyond any fixed $G_D$.  What the graph supplies is a finite, classifiable combinatorial object: adaptive two-adic refinement produces a recurrent core, pair-indexed valuation preimages, a top-bit symmetry, and binomial distance layers.

\end{document}
